% !TEX root = ../../ctfp-print.tex

\lettrine[lhang=0.17]{関}{手について、}圏の構造を保つ写像として話してきた。

関手は一方の圏をもう一方の圏に「埋め込む」。複数のものを1つにまとめることはあっても、
つながりを壊すことはない。別の言い方をすると、関手を使って一方の圏をもう一方の圏の中でモデル化している。
始域圏は、目標の圏の一部となる構造のモデル、設計図として機能する。

\begin{figure}[H]
  \centering\includegraphics[width=0.4\textwidth]{images/1_functors.jpg}
\end{figure}

\noindent
一方の圏をもう一方の圏に埋め込む方法は多数ある場合がある。等価な場合もあれば、
まったく異なる場合もある。始域圏全体を1つの対象に押しつぶすものもあれば、
すべての対象を異なる対象に、すべての射を異なる射に写すものもある。
同じ設計図でも多くの異なる形で実現できる。自然変換はそれらの実現を比較する手段を与えてくれる。
自然変換とは関手間の写像――関手的な性質を保つ特別な写像だ。

圏$\cat{C}$と$\cat{D}$の間の2つの関手$F$と$G$を考えよう。$\cat{C}$の1つの対象$a$に注目すると、
それは$F a$と$G a$という2つの対象に写される。関手の写像は、$F a$を$G a$に写すべきだ。

\begin{figure}[H]
  \centering
  \includegraphics[width=0.3\textwidth]{images/2_natcomp.jpg}
\end{figure}

\noindent
$F a$と$G a$はともに圏$\cat{D}$の対象だ。同じ圏の対象間の写像は圏の流れに逆らうべきではない。
対象間に人為的なつながりを作りたくはない。だから既存のつながり、すなわち射を使うのが\emph{自然}だ。
自然変換とは射の選択だ：すべての対象$a$に対して、$F a$から$G a$への射を1つ選ぶ。
自然変換を$\alpha$と呼ぶと、この射は$a$における$\alpha$の\newterm{成分}と呼ばれ、$\alpha_a$と書く。

\[\alpha_a \Colon F a \to G a\]
$a$は$\cat{C}$における対象であり、$\alpha_a$は$\cat{D}$における射であることを覚えておこう。

ある$a$について、$\cat{D}$に$F a$から$G a$への射が存在しない場合、
$F$と$G$の間の自然変換は存在しない。

もちろんこれは話の半分に過ぎない。関手は対象だけでなく射も写すからだ。では自然変換はそれらの写像に何をするのか？
実は射の写像は固定されている――$F$と$G$の間の任意の自然変換のもとで、$F f$は$G f$に変換されなければならない。
さらに、2つの関手による射の写像は、それと整合する自然変換を定義する際の選択肢を大きく制限する。
$\cat{C}$の2つの対象$a$と$b$の間の射$f$を考えよう。それは$\cat{D}$における2つの射$F f$と$G f$に写される：

\begin{gather*}
  F f \Colon F a \to F b \\
  G f \Colon G a \to G b
\end{gather*}
自然変換$\alpha$は、$\cat{D}$の図式を完成させる2つの射を加える：

\begin{gather*}
  \alpha_a \Colon F a \to G a \\
  \alpha_b \Colon F b \to G b
\end{gather*}

\begin{figure}[H]
  \centering
  \includegraphics[width=0.4\textwidth]{images/3_naturality.jpg}
\end{figure}

\noindent
今や$F a$から$G b$へ向かう2つの経路がある。それらが等しいことを保証するために、
任意の$f$について成り立つ\newterm{自然性条件}を課さなければならない：

\[G f \circ \alpha_a = \alpha_b \circ F f\]
自然性条件はかなり厳しい要件だ。たとえば射$F f$が可逆であれば、
自然性は$\alpha_a$を通じて$\alpha_b$を決定する。つまり$\alpha_a$を$f$に沿って\emph{輸送する}：

\[\alpha_b = (G f) \circ \alpha_a \circ (F f)^{-1}\]

\begin{figure}[H]
  \centering
  \includegraphics[width=0.4\textwidth]{images/4_transport.jpg}
\end{figure}

\noindent
2つの対象の間に複数の可逆な射がある場合、これらの輸送はすべて一致しなければならない。
一般に射は可逆ではないが、2つの関手の間に自然変換が存在するかどうかはまったく保証されないことがわかる。
したがって自然変換で結ばれる関手の少なさや多さは、それらが作用する圏の構造について多くを語ることがある。
極限とYoneda補題の話をするときに、その例をいくつか見ることになる。

成分ごとに見ると、自然変換は対象を射に写すと言える。自然性条件のため、
$\cat{C}$のすべての射に対して$\cat{D}$に自然性の可換正方形が1つあることから、
射を可換正方形に写すとも言える。

\begin{figure}[H]
  \centering
  \includegraphics[width=0.4\textwidth]{images/naturality.jpg}
\end{figure}

\noindent
自然変換のこの性質は、可換図式を多用する圏論的構成において非常に役立つ。
関手を適切に選べば、可換性の条件の多くを自然性の条件に変換できる。
極限、余極限、随伴の話をするときにその例を見ることになる。

最後に、自然変換は関手の同型を定義するために使える。2つの関手が自然同型であるとは、
それらがほぼ同じであるということだ。\newterm{自然同型}は、すべての成分が同型射（可逆な射）である
自然変換として定義される。

\section{多相関数}

プログラミングにおける関手（より具体的にはエンドファンクタ）の役割について話してきた。
それらは型を型に写す型コンストラクタに対応する。また関数を関数に写し、
その写像はC++における高階関数\code{fmap}（または\code{transform}、\code{then}など）として実装される。

自然変換を構成するために、対象、ここでは型\code{a}から始める。一方の関手\code{F}はそれを型$F a$に写す。
もう一方の関手\code{G}はそれを$G a$に写す。\code{a}における自然変換\code{alpha}の成分は
$F a$から$G a$への関数だ。擬似Haskellで書くと：

\begin{snipv}
alpha\textsubscript{a} :: F a -> G a
\end{snipv}
自然変換はすべての型\code{a}に対して定義される多相関数だ：

\src{snippet01}
Haskellでは\code{forall a}は省略可能だ（実際には言語拡張\code{ExplicitForAll}を有効にする必要がある）。
通常はこのように書く：

\src{snippet02}
これは実際には\code{a}によってパラメータ化された関数族であることを覚えておこう。
これはHaskellの構文の簡潔さのもう1つの例だ。C++での同様の構成は少し冗長になる：

\begin{snip}{cpp}
template<typename A> G<A> alpha(F<A>);
\end{snip}
Haskellの多相関数とC++のジェネリック関数には、それらの実装と型チェックの方法に反映される
より深い違いがある。Haskellでは多相関数はすべての型に対して一様に定義されなければならない。
1つの式がすべての型で機能しなければならない。これは\newterm{パラメトリック多相}と呼ばれる。

一方C++はデフォルトで\newterm{アドホック多相}をサポートしており、テンプレートがすべての型に対して
適切に定義されている必要はない。テンプレートが特定の型で機能するかどうかは、
具体的な型が型パラメータに代入されるインスタンス化の時点で決まる。
型チェックが遅延されるため、残念ながら理解しにくいエラーメッセージが生じることが多い。

C++には関数のオーバーロードとテンプレート特殊化の仕組みもあり、異なる型に対して
同じ関数の異なる定義を持つことができる。Haskellではこの機能は型クラスと型族によって提供される。

Haskellのパラメトリック多相には予期しない結果がある：型：

\src{snippet03}
の任意の多相関数は、ここで\code{F}と\code{G}が関手であれば、自動的に自然性条件を満たす。
圏論的記法で書くと（$f$は関数$f \Colon a \to b$）：

\[G f \circ \alpha_a = \alpha_b \circ F f\]
Haskellでは関手\code{G}の射\code{f}への作用は\code{fmap}を使って実装される。
まず明示的な型注釈を使った擬似Haskellで書こう：

\begin{snipv}
fmap\textsubscript{G} f . alpha\textsubscript{a} = alpha\textsubscript{b} . fmap\textsubscript{F} f
\end{snipv}
型推論のおかげでこれらの注釈は不要であり、次の等式が成り立つ：

\begin{snip}{text}
fmap f . alpha = alpha . fmap f
\end{snip}
これはまだ本物のHaskellではない――関数の等式はコードで表現できない――しかし
プログラマが等式推論に使えるか、コンパイラが最適化を実装するために使える恒等式だ。

Haskellで自然性条件が自動的に成り立つ理由は「タダの定理」に関係している。
Haskellで自然変換を定義するために使われるパラメトリック多相は、実装に非常に強い制限を課す――
すべての型に対して1つの式。これらの制限はそのような関数についての等式定理に変換される。
関手を変換する関数の場合、タダの定理が自然性条件となる。\footnote{
  タダの定理についてのより詳しい内容は私のブログ
  \href{https://bartoszmilewski.com/2014/09/22/parametricity-money-for-nothing-and-theorems-for-free/}{「パラメトリシティ：
    タダのお金と定理のための理論」}を参照。}

以前に述べたHaskellの関手の考え方の1つは、一般化されたコンテナとして考えることだ。
この類推を続けると、自然変換を1つのコンテナの内容を別のコンテナに詰め替えるレシピと見なすことができる。
アイテム自体には触れない：変更もしないし、新しいものも作らない。
ただ（一部の）アイテムを新しいコンテナにコピーし、時に複数回コピーするだけだ。

自然性条件は次のことを述べているとも言える：先に\code{fmap}の適用によってアイテムを変更してから
後で詰め替えることと、先に詰め替えてから新しいコンテナで\code{fmap}の独自実装によってアイテムを
変更することは、どちらでも同じ結果になる。詰め替えと\code{fmap}適用という2つの操作は直交している。
「一方は卵を移動し、他方は卵を茹でる。」

Haskellでの自然変換の例をいくつか見てみよう。最初はリスト関手と\code{Maybe}関手の間のものだ。
リストの先頭を返すが、リストが空でない場合のみ：

\src{snippet04}
これは\code{a}において多相な関数だ。制限なしにすべての型\code{a}で機能するので、
パラメトリック多相の例だ。したがって2つの関手間の自然変換だ。しかし確認のために、自然性条件を検証してみよう。

\src{snippet05}
考慮すべき2つのケースがある：空リストの場合：

\src{snippet06}

\src{snippet07}
非空リストの場合：

\src{snippet08}

\src{snippet09}
リストに対する\code{fmap}の実装を使った：

\src{snippet10}
\code{Maybe}に対しては：

\src{snippet11}
興味深いケースは関手の1つが自明な\code{Const}関手のときだ。\code{Const}関手との間の
自然変換は、戻り値の型または引数の型において多相な関数のように見える。

たとえば\code{length}はリスト関手から\code{Const Int}関手への自然変換と見なせる：

\src{snippet12}
ここで\code{unConst}は\code{Const}コンストラクタを取り除くために使われる：

\src{snippet13}
もちろん実際には\code{length}はこのように定義される：

\src{snippet14}
これは実際には自然変換であるという事実を効果的に隠している。

\code{Const}関手\emph{から}のパラメトリック多相関数を見つけるのはやや難しい。
なぜなら無から値を作る必要があるからだ。できることは：

\src{snippet15}
すでに見たもう1つの重要な関手、Yoneda補題で重要な役割を果たす\code{Reader}関手だ。
その定義を\code{newtype}として書き直そう：

\src{snippet16}
これは2つの型でパラメータ化されているが、（共変的に）関手的なのは2番目の型についてのみだ：

\src{snippet17}
すべての型\code{e}に対して、\code{Reader e}から他の任意の関手\code{f}への自然変換の族を定義できる。
後でこの族のメンバーが常に\code{f e}の元と1対1対応していることを見る
（\hyperref[the-yoneda-lemma]{Yoneda補題}）。

たとえば、1つの元\code{()}を持つ多少自明な単位型\code{()}を考えよう。
関手\code{Reader ()}は任意の型\code{a}を関数型\code{() -> a}に写す。
これらは集合\code{a}から単一の元を選ぶすべての関数だ。これらの数は\code{a}の元の数と同じだ。
次にこの関手から\code{Maybe}関手への自然変換を考えよう：

\src{snippet18}
これらは\code{dumb}と\code{obvious}の2つだけだ：

\src{snippet19}
そして：

\src{snippet20}
（\code{g}でできることは、単位値\code{()}に適用することだけだ。）

そして確かに、Yoneda補題が予言するように、これらは\code{Nothing}と\code{Just ()}である
\code{Maybe ()}型の2つの元に対応する。Yoneda補題については後で再び取り上げる――
これはちょっとしたティーザーだった。

\section{自然性を超えて}

2つの関手（\code{Const}関手の端点ケースを含む）の間のパラメトリック多相関数は
常に自然変換だ。すべての標準的な代数的データ型は関手なので、そのような型間の任意の多相関数は
自然変換だ。

また関数型も使えるが、それらは戻り値の型において関手的だ。関手（\code{Reader}関手のような）を
構築し、高階関数である自然変換を定義するために使える。

しかし関数型は引数の型において共変ではない。それらは\newterm{反変}だ。
もちろん反変関手は反対圏からの共変関手と等価だ。2つの反変関手間の多相関数は
圏論的な意味でまだ自然変換だが、それらは反対圏からHaskellの型への関手について機能する。

以前に見た反変関手の例を覚えているかもしれない：

\src{snippet21}
この関手は\code{a}において反変だ：

\src{snippet22}
たとえば\code{Op Bool}から\code{Op String}への多相関数を書けるが：

\src{snippet23}
2つの関手は共変ではないため、これは$\Hask$における自然変換ではない。
しかし両方とも反変なので、「反対の」自然性条件を満たす：

\src{snippet24}[b]
\code{contramap}のシグネチャのため、関数\code{f}は\code{fmap}で使うものとは
反対方向に進まなければならないことに注意：

\src{snippet25}
共変または反変のいずれでもない型コンストラクタはあるか？例を挙げよう：

\src{snippet26}
同じ型\code{a}が負の位置（反変）と正の位置（共変）の両方で使われているため、
これは関手ではない。この型に対して\code{fmap}も\code{contramap}も実装できない。
したがってシグネチャ：

\src{snippet27}
の関数は、\code{f}が任意の関手のとき、自然変換にはなれない。
興味深いことに、そのようなケースを扱う自然変換の一般化、二重自然変換と呼ばれるものがある。
エンドを議論するときにそれらについて取り上げる。

\section{関手圏}

関手間の写像――自然変換――ができたので、関手が圏を形成するかどうかを問うのは自然なことだ。
そして確かに形成する！圏の各ペア$\cat{C}$と$\cat{D}$に対して、関手の圏が1つある。
この圏の対象は$\cat{C}$から$\cat{D}$への関手であり、射はそれらの関手間の自然変換だ。

2つの自然変換の合成を定義する必要があるが、これはかなり簡単だ。自然変換の成分は射であり、
射の合成の仕方は知っている。

実際、関手$F$から$G$への自然変換$\alpha$を取ろう。対象$a$におけるその成分はある射だ：
\[\alpha_a \Colon F a \to G a\]
$\alpha$と、関手$G$から$H$への自然変換$\beta$を合成したい。$a$における$\beta$の成分は射だ：
\[\beta_a \Colon G a \to H a\]
これらの射は合成可能で、その合成は別の射となる：
\[\beta_a \circ \alpha_a \Colon F a \to H a\]
この射を、$\beta$の後の$\alpha$という2つの自然変換の合成である自然変換$\beta \cdot \alpha$の成分として使う：
\[(\beta \cdot \alpha)_a = \beta_a \circ \alpha_a\]

\begin{figure}[H]
  \centering
  \includegraphics[width=0.4\textwidth]{images/5_vertical.jpg}
\end{figure}

\noindent
図式を（長い間）見れば、この合成の結果が確かにFからHへの自然変換であることがわかる：
\[H f \circ (\beta \cdot \alpha)_a = (\beta \cdot \alpha)_b \circ F f\]

\begin{figure}[H]
  \centering
  \includegraphics[width=0.35\textwidth]{images/6_verticalnaturality.jpg}
\end{figure}

\noindent
自然変換の合成は結合的だ。なぜなら成分である通常の射は合成について結合的だからだ。

最後に、各関手$F$に対して恒等自然変換$1_F$が存在し、その成分は恒等射だ：
\[\id_{F a} \Colon F a \to F a\]
したがって確かに、関手は圏を形成する。

記法について一言。Saunders Mac Laneに倣い、今述べた種類の自然変換の合成にドットを使う。
問題は自然変換を合成する方法が2つあることだ。これは垂直合成と呼ばれる。
なぜなら関手はそれを記述する図式で通常縦に積み重ねられるからだ。
垂直合成は関手圏を定義するうえで重要だ。水平合成についてはすぐに説明する。

\begin{figure}[H]
  \centering
  \includegraphics[width=0.3\textwidth]{images/6a_vertical.jpg}
\end{figure}

\noindent
圏$\cat{C}$と$\cat{D}$の間の関手圏は$\cat{Fun(C, D)}$または$\cat{{[}C, D{]}}$、
あるいは$\cat{D^C}$と書く。この最後の記法は、関手圏自体が他のある圏における
関数対象（指数）と見なせることを示唆している。実際そうなのか？

これまで構築してきた抽象の階層を見てみよう。対象と射の集まりである圏から始めた。
圏自体（より厳密には、対象が集合をなす\emph{小さい}圏）は高いレベルの圏$\Cat$における
対象だ。その圏における射は関手だ。$\Cat$のホム集合は関手の集合だ。たとえば
$\cat{Cat(C, D)}$は2つの圏$\cat{C}$と$\cat{D}$の間の関手の集合だ。

\begin{figure}[H]
  \centering
  \includegraphics[width=0.3\textwidth]{images/7_cathomset.jpg}
\end{figure}

\noindent
関手圏$\cat{{[}C, D{]}}$も2つの圏の間の関手の集合（さらに射としての自然変換）だ。
その対象は$\cat{Cat(C, D)}$のメンバーと同じだ。さらに関手圏は圏として$\Cat$の対象でなければならない
（2つの小さい圏の間の関手圏はそれ自体小さい）。圏のホム集合と同じ圏の対象の間の関係がある。
この状況は前のセクションで見た指数対象とまったく同じだ。$\Cat$においてそれをどう構成するか見てみよう。

覚えているように、指数を構成するためにはまず積を定義する必要がある。$\Cat$では、
小さい圏は対象の\emph{集合}なので、集合のデカルト積を定義する方法を知っているため、これは比較的簡単だ。
したがって積圏$\cat{C\times D}$の対象は単に対象のペア$(c, d)$で、一方は$\cat{C}$から、
もう一方は$\cat{D}$からだ。同様に、2つのそのようなペア$(c, d)$と$(c', d')$の間の射は
射のペア$(f, g)$で、$f \Colon c \to c'$かつ$g \Colon d \to d'$だ。
これらの射のペアは成分ごとに合成され、常に恒等射のペアである恒等ペアが存在する。
簡単に言えば、$\Cat$は任意の圏のペアに対して指数対象$\cat{D^C}$を持つ
完全なデカルト閉圏だ。そして$\Cat$における「対象」とは圏のことなので、
$\cat{D^C}$は圏であり、$\cat{C}$と$\cat{D}$の間の関手圏として同定できる。

\section{2圏}

それはさておき、$\Cat$をより詳しく見てみよう。定義により$\Cat$の任意のホム集合は
関手の集合だ。しかし見てきたように、2つの対象間の関手は単なる集合よりも豊かな構造を持つ。
それらは圏を形成し、射として自然変換が作用する。関手は$\Cat$における射と見なされるので、
自然変換は射間の射だ。

この豊かな構造は$\cat{2}$圏の例で、対象と射（この文脈では$1$射と呼ばれることもある）に加えて、
射間の射である$2$射も存在する圏の一般化だ。

$\Cat$を$\cat{2}$圏として見た場合：

\begin{itemize}
  \tightlist
  \item
        対象：（小さい）圏
  \item
        1射：圏間の関手
  \item
        2射：関手間の自然変換
\end{itemize}

\begin{figure}[H]
  \centering
  \includegraphics[width=0.3\textwidth]{images/8_cat-2-cat.jpg}
\end{figure}

\noindent
2つの圏$\cat{C}$と$\cat{D}$の間のホム集合の代わりに、ホム圏――関手圏$\cat{D^C}$――がある。
通常の関手の合成がある：$\cat{D^C}$からの関手$F$は$\cat{E^D}$からの関手$G$と合成されて
$\cat{E^C}$からの$G \circ F$を与える。しかし各ホム圏の内部にも合成がある――
関手間の自然変換、すなわち2射の垂直合成だ。

$\cat{2}$圏に2種類の合成があると、問題が生じる：それらはお互いにどのように作用するのか？

$\Cat$の2つの関手、すなわち1射を選ぼう：
\begin{gather*}
  F \Colon \cat{C} \to \cat{D} \\
  G \Colon \cat{D} \to \cat{E}
\end{gather*}
とその合成：
\[G \circ F \Colon \cat{C} \to \cat{E}\]
それぞれ関手$F$と$G$に作用する2つの自然変換$\alpha$と$\beta$があるとする：
\begin{gather*}
  \alpha \Colon F \to F' \\
  \beta \Colon G \to G'
\end{gather*}

\begin{figure}[H]
  \centering
  \includegraphics[width=0.4\textwidth]{images/10_horizontal.jpg}
\end{figure}

\noindent
$\alpha$の終域が$\beta$の始域と異なるため、このペアに垂直合成を適用できないことに注意。
実際にそれらは2つの異なる関手圏のメンバーだ：$\cat{D^C}$と$\cat{E^D}$。
しかし関手$F'$と$G'$に合成を適用できる。なぜなら$F'$の終域が$G'$の始域――圏$\cat{D}$――と
一致するからだ。では関手$G' \circ F'$と$G \circ F$の関係は何か？

$\alpha$と$\beta$を使って、$G \circ F$から$G' \circ F'$への自然変換を定義できるか？
構成の概略を示そう。

\begin{figure}[H]
  \centering
  \includegraphics[width=0.5\textwidth]{images/9_horizontal.jpg}
\end{figure}

\noindent
いつものように$\cat{C}$の対象$a$から始める。その像は$\cat{D}$の2つの対象$F a$と$F'a$に分かれる。
これら2つの対象を結ぶ$\alpha$の成分である射もある：
\[\alpha_a \Colon F a \to F'a\]
$\cat{D}$から$\cat{E}$へ向かうと、これらの2つの対象はさらに4つの対象$G (F a)$、$G'(F a)$、
$G (F'a)$、$G'(F'a)$に分かれる。また正方形を形成する4つの射もある。
これらのうち2つは自然変換$\beta$の成分だ：
\begin{gather*}
  \beta_{F a} \Colon G (F a) \to G'(F a) \\
  \beta_{F'a} \Colon G (F'a) \to G'(F'a)
\end{gather*}
他の2つは2つの関手による$\alpha_a$の像だ（関手は射を写す）：
\begin{gather*}
  G \alpha_a \Colon G (F a) \to G (F'a) \\
  G'\alpha_a \Colon G'(F a) \to G'(F'a)
\end{gather*}
多くの射がある。目標は$G (F a)$から$G'(F'a)$への射を見つけることで、
これが2つの関手$G \circ F$と$G' \circ F'$を結ぶ自然変換の成分の候補となる。
実は$G (F a)$から$G'(F'a)$へ向かう経路は1つではなく2つある：
\begin{gather*}
  G'\alpha_a \circ \beta_{F a} \\
  \beta_{F'a} \circ G \alpha_a
\end{gather*}
幸運なことに、それらは等しい。なぜなら形成した正方形が$\beta$の自然性正方形であることがわかるからだ。

$G \circ F$から$G' \circ F'$への自然変換の成分を定義した。この変換の自然性の証明は、
十分な根気があればかなり直接的だ。

この自然変換を$\alpha$と$\beta$の\newterm{水平合成}と呼ぶ：
\[\beta \circ \alpha \Colon G \circ F \to G' \circ F'\]
再びMac Laneに倣い、水平合成には小さな丸を使う。ただし代わりにアスタリスクを使うこともある。

圏論的な経験則がある：合成があるときはいつも圏を探せ。垂直合成の自然変換があり、
それは関手圏の一部だ。では水平合成はどうか？それはどの圏に属するのか？

これを解明する方法は$\Cat$を横向きに見ることだ。自然変換を関手間の矢印としてではなく、
圏間の矢印として見る。自然変換は2つの圏の間に位置し、変換する関手で結ばれている。
それらの2つの圏を結んでいると考えることができる。

\begin{figure}[H]
  \centering
  \includegraphics[width=0.5\textwidth]{images/sideways.jpg}
\end{figure}

\noindent
$\Cat$の2つの対象――圏$\cat{C}$と$\cat{D}$――に注目しよう。$\cat{C}$と$\cat{D}$を結ぶ関手の間を
行き来する自然変換の集合がある。これらの自然変換が$\cat{C}$から$\cat{D}$への新しい矢印となる。
同様に、$\cat{D}$から$\cat{E}$を結ぶ関手の間の自然変換があり、それらを$\cat{D}$から$\cat{E}$への
新しい矢印として扱える。水平合成はこれらの矢印の合成だ。

また$\cat{C}$から$\cat{C}$への恒等矢印もある。それは$\cat{C}$の恒等関手をそれ自体に写す
恒等自然変換だ。水平合成の恒等は垂直合成の恒等でもあるが、その逆は成り立たないことに注意。

最後に、2つの合成は交換法則を満たす：
\[(\beta' \cdot \alpha') \circ (\beta \cdot \alpha) = (\beta' \circ \beta) \cdot (\alpha' \circ \alpha)\]
ここでSaunders Mac Laneの言葉を引用しよう：「読者はこれを証明するのに必要な明らかな図式を
書き出すことを楽しむかもしれない。」

将来役立つかもしれないもう1つの記法がある。この$\Cat$の新しい横向きの解釈では、
対象から対象へ向かう方法が2つある：関手を使うか自然変換を使うかだ。
しかし関手の矢印を特別な種類の自然変換として再解釈できる：この関手に作用する恒等自然変換として。
だからこの記法をよく見かけることになる：
\[F \circ \alpha\]
ここで$F$は$\cat{D}$から$\cat{E}$への関手であり、$\alpha$は$\cat{C}$から$\cat{D}$への
2つの関手の間の自然変換だ。関手と自然変換を合成できないので、これは$\alpha$の後の恒等自然変換
$1_F$の水平合成として解釈される。

同様に：
\[\alpha \circ F\]
は$1_F$の後の$\alpha$の水平合成だ。

\section{まとめ}

これで本書の第1部が終了した。圏論の基本的な語彙を学んだ。対象と圏を名詞として、
射、関手、自然変換を動詞として考えることができる。射は対象をつなぎ、関手は圏をつなぎ、
自然変換は関手をつなぐ。

しかし同時に、ある抽象レベルで行為として現れるものが、次のレベルでは対象となることも見てきた。
射の集合が関数対象になる。対象として、別の射の始域や終域になれる。これが高階関数の背後にある考え方だ。

関手は対象を対象に写すので、型コンストラクタ、すなわちパラメトリック型として使える。
関手はまた射を写すので高階関数だ――\code{fmap}。\code{Const}、積、余積のような
単純な関手があり、それらは多種多様な代数的データ型を生成するために使える。
関数型も関手的で、共変と反変の両方があり、代数的データ型を拡張するために使える。

関手は関手圏の対象と見なすことができる。そうすることで射の始域と終域になる：自然変換が。
自然変換は多相関数の特別な型だ。

\section{演習問題}

\begin{enumerate}
  \tightlist
  \item
        \code{Maybe}関手からリスト関手への自然変換を定義せよ。そのための自然性条件を証明せよ。
  \item
        \code{Reader ()}とリスト関手の間の少なくとも2つの異なる自然変換を定義せよ。
        \code{()}のリストはいくつあるか？
  \item
        前の演習を\code{Reader Bool}と\code{Maybe}で続けよ。
  \item
        自然変換の水平合成が自然性条件を満たすことを示せ（ヒント：成分を使え）。
        図式の追跡のよい演習だ。
  \item
        交換法則を証明するのに必要な明らかな図式を書き出すことをどのように楽しめるかについて
        短いエッセイを書け。
  \item
        異なる\code{Op}関手間の変換の反対の自然性条件のいくつかのテストケースを作れ。
        次の選択肢がある：

        \begin{snip}{haskell}
op :: Op Bool Int
op = Op (\x -> x > 0)
\end{snip}
        そして：

        \begin{snip}{haskell}
f :: String -> Int
f x = read x
\end{snip}
\end{enumerate}
